% Lecture 3 — Vector spaces
% Compiled to web slides by tools/tex2slides.mjs

\lecture{Vector spaces}{Linear Algebra · Warsaw University of Technology}{Lecture 3 · Tomasz Brengos}

\section{Motivation}

\begin{frame}{Turn your data into a vector}
If there is one habit that defines modern computing, it is this: \emph{turn your data into a vector}.

\begin{itemize}
\item A $1000\times 1000$ greyscale image is a list of a million pixel brightnesses — a vector in $\R^{1\,000\,000}$.
\pause
\item An audio clip is a vector of samples; a streaming user is a vector of ratings.
\pause
\item A word, inside a language model, is a vector of a few hundred learned coordinates.
\end{itemize}

\pause

\begin{remark}
Once data is a vector, geometry becomes computation: similar users are \emph{nearby} vectors, and
recommendation is finding neighbours.
\end{remark}
\note{Sell the abstraction: this chapter axiomatises exactly the operations used to blend images and mix audio.}
\end{frame}

\begin{frame}{Vector arithmetic captures meaning}
Adding vectors and scaling them is how images are blended and audio tracks are mixed. More strikingly, in
word-vector models the sum and difference of vectors captures \emph{meaning}:
\[
\text{king} - \text{man} + \text{woman} \approx \text{queen}
\]
is literally vector arithmetic in the sense of this chapter.

\pause

\begin{remark}
We will not define vectors as ``arrows'' or ``tuples'', but as elements of \emph{any} set with a
well-behaved addition and scaling. The payoff: tuples, polynomials, matrices and functions all satisfy the
same axioms, so every theorem below holds for all of them at once.
\end{remark}
\end{frame}

\begin{frame}{The question this chapter answers}
How many numbers does it really take to describe your data?

\begin{itemize}
\item A \emph{spanning set} is a collection of building blocks capable of expressing everything in sight.
\pause
\item \emph{Linear independence} means no building block is redundant.
\pause
\item A \emph{basis} is a minimal, lossless set of building blocks; its size — the \emph{dimension} — is
the honest number of degrees of freedom of the data.
\end{itemize}

\pause

\begin{remark}
Real data rarely uses all of its apparent dimensions: photographs of faces occupy a tiny, well-approximated
sliver of $\R^{1\,000\,000}$ — the reason image compression and recommender systems work at all.
\end{remark}
\end{frame}

\section{Vector spaces}

\begin{frame}{Where the axioms come from}
You already know $\R^2 = \{(x,y)\mid x,y\in\R\}$, where vectors add and scale:
\[
(x_1,y_1)\oplus(x_2,y_2) = (x_1+x_2,\, y_1+y_2), \qquad \lambda\odot(x,y) = (\lambda x,\, \lambda y).
\]

\pause

The abstract vector space axiomatises exactly these two operations, for \emph{any} set and \emph{any} field
of scalars.
\end{frame}

\begin{frame}{Definition}
\begin{definition}
Let $\K$ be a field. A set $\V$ with vector addition $+$ and scalar multiplication $\cdot\colon\K\times\V\to\V$
is a \emph{vector space} over $\K$ if:
\begin{itemize}
\item $\uu+\vv=\vv+\uu$ and $\uu+(\vv+\ww)=(\uu+\vv)+\ww$;
\item there is $\0\in\V$ with $\vv+\0=\vv$ for all $\vv$;
\item every $\vv$ has $-\vv$ with $\vv+(-\vv)=\0$;
\item $\lambda\cdot(\uu+\vv)=\lambda\cdot\uu+\lambda\cdot\vv$ and $(\lambda_1+\lambda_2)\cdot\uu=\lambda_1\cdot\uu+\lambda_2\cdot\uu$;
\item $(\lambda_1\lambda_2)\cdot\uu=\lambda_1\cdot(\lambda_2\cdot\uu)$ and $1\cdot\uu=\uu$.
\end{itemize}
\end{definition}
\note{Don't dwell on each axiom individually — the point is that this is exactly the R^2 arithmetic, generalised.}
\end{frame}

\begin{frame}{Free consequences of the axioms}
\begin{fact}
For any $\vv\in\V$ and $\lambda\in\K$:
\[
0\cdot\vv=\0, \qquad \lambda\cdot\0=\0, \qquad (-1)\cdot\vv=-\vv, \qquad \lambda\cdot\vv=\0 \implies \lambda=0 \text{ or } \vv=\0.
\]
\end{fact}

\pause

\begin{steps}
\step[distributivity]{0\cdot \vv = (0+0)\cdot \vv = 0\cdot\vv + 0\cdot \vv}
\step[add $-(0\cdot\vv)$ to both sides]{\0 = 0\cdot\vv}
\end{steps}
\end{frame}

\begin{frame}{Example: K to the n}
\begin{example}
For any field $\K$ and $n\in\N$, the set $\K^n=\{(x_1,\ldots,x_n)\mid x_i\in\K\}$ with coordinatewise
addition and scaling is a vector space, with zero vector $\0=(0,\ldots,0)$.

Taking $\K=\R,n=2$ recovers $\R^2$. Taking $\K=\C$ gives $\C^n$.
\end{example}
\end{frame}

\begin{frame}{Example: bit vectors}
\begin{example}
Taking $\K=\Z_2$ gives the space $\Z_2^n$ of \emph{bit vectors}: addition is bitwise XOR, e.g.\ in $\Z_2^4$,
\[
(1,0,1,1)+(1,1,0,1) = (0,1,1,0).
\]

\pause

Every vector is its own additive inverse. Identifying a subset of $\{1,\ldots,n\}$ with its characteristic
vector, $\Z_2^n$ is also the set of all subsets of $\{1,\ldots,n\}$, addition becoming symmetric difference.
\end{example}

\pause

\begin{remark}
Vector spaces over $\Z_2$ are the natural home of error-correcting codes and appear throughout cryptography.
\end{remark}
\end{frame}

\begin{frame}{More examples}
\begin{columns}
\column
\begin{example}
$\C$ over $\R$: standard addition and multiplication by real scalars makes $\C$ a real vector space.
\end{example}
\column
\begin{example}
Polynomials $\K[x]$, with the standard operations, form a vector space; so does $\K_n[x]$, the polynomials
of degree $\leq n$.
\end{example}
\end{columns}

\pause

\begin{example}
For a non-empty set $X$, the functions $\K^X = \{f\colon X\to\K\}$ with pointwise operations
$(f+g)(x)=f(x)+g(x)$, $(\lambda\cdot f)(x)=\lambda f(x)$ form a vector space.
\end{example}
\end{frame}

\section{Subspaces}

\begin{frame}{Definition}
\begin{definition}
$\W\subseteq\V$ is a \emph{subspace} of $\V$ if $\W$ is itself a vector space under the same operations.
\end{definition}

\pause

\begin{remark}[A shortcut]
It suffices to check three things: $\0\in\W$ (in particular $\W\neq\emptyset$), and $\W$ is closed under
$+$ and under scalar multiplication — since $\0=0\cdot\ww$ and $-\ww=(-1)\cdot\ww$ come for free. Every
$\V$ has the trivial subspaces $\{\0\}$ and $\V$.
\end{remark}
\end{frame}

\begin{frame}{A plane through the origin}
\begin{example}
In $\R^3$, $W=\{(x,y,z)\mid 2x+y+3z=0\}$ is a subspace: $\0\in W$, and if $\ww_1,\ww_2\in W$ then
\[
2(x_1+x_2)+(y_1+y_2)+3(z_1+z_2) = \underbrace{(2x_1+y_1+3z_1)}_{=0} + \underbrace{(2x_2+y_2+3z_2)}_{=0} = 0,
\]
so $W$ is closed under addition (and similarly under scaling).
\end{example}

\pause

\begin{warning}
$W' = \{(x,y,z)\mid x^2=1\}$ is \emph{not} a subspace: $\0\notin W'$, and $(1,0,0)+(1,0,0)=(2,0,0)\notin W'$.
\end{warning}
\end{frame}

\begin{frame}{The union trap}
\begin{warning}
The \emph{intersection} of two subspaces is always a subspace — but their \emph{union} usually is not.
\end{warning}

\pause

In $\R^2$ the axes $W_1=\{(x,0)\}$ and $W_2=\{(0,y)\}$ are subspaces, yet
\[
(1,0)+(0,1) = (1,1) \notin W_1\cup W_2.
\]
\end{frame}

\section{Span}

\begin{frame}{Linear combinations}
\begin{definition}
For scalars $a_1,\ldots,a_n\in\K$ and vectors $\vv_1,\ldots,\vv_n\in\V$, the vector
$a_1\vv_1+\ldots+a_n\vv_n$ is their \emph{linear combination}.
\end{definition}

\begin{example}
In $\R^2$: $(2,3) = 2\cdot(1,1)+1\cdot(0,1)$.
\end{example}
\end{frame}

\begin{frame}{Span, and playing with it}
\begin{columns}
\column
\begin{definition}
$\Span{S} = \{a_1\vv_1+\ldots+a_n\vv_n \mid a_i\in\K,\ \vv_i\in S\}$ — all linear combinations of vectors
from $S$. By convention $\Span{\emptyset}=\{\0\}$.
\end{definition}

Drag $u$ and $v$, then slide $a,b$: the marked point is $au+bv$, ranging over $\Span{\{u,v\}}$.
\column
\widget{VecCombo}{}
\end{columns}
\end{frame}

\begin{frame}{Span is a subspace}
\begin{theorem}
$\Span{\{\vv_1,\ldots,\vv_n\}}$ is a subspace of $\V$ containing $\vv_1,\ldots,\vv_n$.
\end{theorem}

\pause

\begin{proof}
For $\uu=\sum a_i\vv_i$, $\ww=\sum b_i\vv_i$: $\uu+\ww=\sum(a_i+b_i)\vv_i$, $-\uu=\sum(-a_i)\vv_i$,
$\0=\sum 0\cdot\vv_i$, and $\lambda\uu=\sum(\lambda a_i)\vv_i$ — all back in the span.
\end{proof}
\end{frame}

\begin{frame}{Properties of span}
\begin{fact}
For $S,T\subseteq\V$:
\begin{itemize}
\item $S\subseteq\Span{S}$;
\pause
\item $S\subseteq T \implies \Span{S}\subseteq\Span{T}$;
\pause
\item $\Span{\Span{S}} = \Span{S}$;
\pause
\item $\Span{S}$ is the \emph{smallest} subspace containing $S$.
\end{itemize}
\end{fact}

\pause

\begin{example}
$\Span{\{1,i\}}=\C$ over $\R$, but $\Span{\{3+i\}} = \{3a+ai\mid a\in\R\}\neq\C$.
\end{example}
\end{frame}

\begin{frame}{Testing membership}
Is $(1,2,3)\in\Span{\{(1,1,1),(1,0,1)\}}$ in $\R^3$? Solve $a(1,1,1)+b(1,0,1)=(1,2,3)$:

\pause

\begin{steps}
\step[the three equations]{a+b=1, \quad a=2, \quad a+b=3}
\step[the first and third contradict]{\text{no solution} \implies (1,2,3)\notin\Span{\ldots}}
\end{steps}

\pause

For $(3,1,3)$ the same equations give $a=1,b=2$ consistently:
\[
(3,1,3) = 1\cdot(1,1,1)+2\cdot(1,0,1)\in\Span{\{(1,1,1),(1,0,1)\}}.
\]
\note{Membership in a span always reduces to solving a linear system — the systematic theory comes in the next chapter.}
\end{frame}

\section{Independence and basis}

\begin{frame}{Linear independence}
\begin{definition}
$S\subseteq\V$ is \emph{linearly independent} if for any $\vv_1,\ldots,\vv_n\in S$,
\[
a_1\vv_1+\ldots+a_n\vv_n=\0 \implies a_1=\ldots=a_n=0.
\]
Otherwise $S$ is \emph{linearly dependent}.
\end{definition}

\pause

\begin{example}
In $\R^2$: $\{(0,1),(1,1)\}$ is independent, since $a(0,1)+b(1,1)=(b,a+b)=(0,0)$ forces $a=b=0$.
\end{example}
\end{frame}

\begin{frame}{More independence examples}
\begin{itemize}
\item $\C$ over $\R$: $\{1,i\}$ is independent — $a+bi=0\implies a=b=0$.
\pause
\item $\R[x]$: $\{x+1,\,x-1,\,2\}$ is \emph{dependent}: $(x+1)-(x-1)-2=0$.
\pause
\item $\R[x]$: $\{1,x,x^2,x^3,\ldots\}$ is independent.
\end{itemize}
\end{frame}

\begin{frame}{Testing independence: a worked example}
Are $(1,2,3),(2,1,0),(1,-4,-9)$ independent in $\R^3$? Solve $a(1,2,3)+b(2,1,0)+c(1,-4,-9)=\0$:

\pause

\begin{steps}
\step[three equations]{a+2b+c=0, \quad 2a+b-4c=0, \quad 3a-9c=0}
\step[third gives a in terms of c]{a=3c}
\step[substitute into the first]{b=-2c}
\step[the second is then automatic; take c=1]{3(1,2,3) - 2(2,1,0) + 1(1,-4,-9) = (0,0,0)}
\end{steps}

\pause

A non-trivial combination equals $\0$: the vectors are \emph{dependent}.
\end{frame}

\begin{frame}{Two useful facts}
\begin{fact}
$\{\vv_1,\ldots,\vv_n\}$ is dependent iff some $\vv_i$ is a linear combination of the others.
\end{fact}

\pause

\begin{fact}
Every subset of an independent set is independent. Equivalently: any set containing a dependent subset is
dependent. In particular, any set containing $\0$ is dependent, since $1\cdot\0=\0$.
\end{fact}
\end{frame}

\begin{frame}{Basis}
\begin{definition}
$B\subseteq\V$ is a \emph{basis} of $\V$ if $B$ is linearly independent and $\Span{B}=\V$.
\end{definition}

\pause

\begin{itemize}
\item $\R^2$: $\{(1,0),(1,1)\}$ is a basis.
\pause
\item $\C$ over $\R$: $\{1,i\}$ is a basis; so is $\{1,i+2\}$.
\pause
\item $\K^n$: the \emph{standard basis} $\{e_1,\ldots,e_n\}$, $e_i=(0,\ldots,1,\ldots,0)$.
\pause
\item $\K_n[x]$: the basis $\{1,x,x^2,\ldots,x^n\}$.
\end{itemize}
\end{frame}

\begin{frame}{Basis = unique representation}
\begin{theorem}
$\{\vv_1,\ldots,\vv_n\}$ is a basis of $\V$ iff every $\ww\in\V$ is expressible \emph{uniquely} as a linear
combination of the $\vv_i$.
\end{theorem}

\pause

\begin{proof}
($\Rightarrow$) Two representations subtract to $\sum(a_i-b_i)\vv_i=\0$; independence forces $a_i=b_i$.\\
($\Leftarrow$) Spanning is immediate. For independence: $\0$ has the representation $\sum 0\cdot\vv_i$;
uniqueness forces any $\sum a_i\vv_i=\0$ to have all $a_i=0$.
\end{proof}
\end{frame}

\begin{frame}{Coordinates}
\begin{definition}
For a basis $B=\{\vv_1,\ldots,\vv_n\}$ (fixed order) and $\ww=a_1\vv_1+\ldots+a_n\vv_n$, the scalars
$(a_1,\ldots,a_n)$ are the \emph{coordinates} of $\ww$ in $B$.
\end{definition}

\pause

\begin{example}
In $\R^2$ with $B=\{(1,0),(1,1)\}$: $(2,3)=-1\cdot(1,0)+3\cdot(1,1)$, so the coordinates are $(-1,3)$ — but
$(2,3)$ itself in the standard basis.
\end{example}

\pause

\begin{remark}
This is the real power of coordinates: once a basis is fixed, \emph{any} vector — polynomial, matrix,
function — is faithfully encoded as a tuple in $\K^n$. All the matrix machinery ahead runs on this bridge.
\end{remark}
\end{frame}

\section{Dimension}

\begin{frame}{Dimension is well defined}
\begin{theorem}
If $\V$ has a finite basis, any two bases of $\V$ have the same size.
\end{theorem}

\pause

\begin{definition}
The \emph{dimension} $\dim(\V)$ is the size of any basis (else $\infty$).
\end{definition}

\pause

\begin{itemize}
\item $\dim(\{\0\})=0$, \quad $\dim(\K^n)=n$, \quad $\dim(\K[x])=\infty$, \quad $\dim(\K_n[x])=n+1$;
\item $\dim(\C \text{ over } \R)=2$, \quad $\dim(\C \text{ over } \C)=1$.
\end{itemize}
\end{frame}

\begin{frame}{Counting theorems}
\begin{theorem}
If $\W\leq\V$: $\dim(\W)\leq\dim(\V)$, and if $\V$ is finite-dimensional, $\dim(\W)=\dim(\V)\implies\W=\V$.
\end{theorem}

\pause

\begin{theorem}
Let $\dim(\V)=n$ and $S\subseteq\V$. Then:
\begin{itemize}
\item $\Span{S}=\V \implies |S|\geq n$;
\item $S$ independent and $|S|=n \implies S$ is a basis;
\item $\Span{S}=\V$ and $|S|=n \implies S$ is a basis;
\item $|S|>n \implies S$ is dependent.
\end{itemize}
\end{theorem}
\end{frame}

\begin{frame}{Extending to a basis}
\begin{theorem}
If $\dim(\V)=n<\infty$: every independent $S\subseteq\V$ extends to a basis, and every spanning
$S\subseteq\V$ contains a basis.
\end{theorem}

\pause

\begin{example}
$\{(1,1,0)\}$ is independent in $\R^3$; it extends to the basis $\{(1,1,0),(0,1,0),(0,0,1)\}$ — three
independent vectors in a $3$-dimensional space automatically form a basis.
\end{example}
\end{frame}

\begin{frame}{Summary}
\begin{itemize}
\item A \emph{vector space} axiomatises addition and scalar multiplication — for any set, any field.
\pause
\item A \emph{subspace}: closed under $+$ and scaling, contains $\0$.
\pause
\item \emph{Span}: all linear combinations; \emph{independence}: no redundancy.
\pause
\item A \emph{basis} is a minimal spanning set; its size is the \emph{dimension}.
\pause
\item \textbf{Next lecture:} matrices — the concrete bridge from abstract vectors to columns of numbers.
\end{itemize}
\end{frame}
